\chapter*{Заключение}
\addcontentsline{toc}{chapter}{Заключение}

В ходе выполнения настоящей преддипломной практики была успешно решена
задача разработки и верификации программно-модельного комплекса для анализа
динамики опухолевых процессов при иммунотерапии. Актуальность исследования
обусловлена необходимостью создания предиктивных инструментов для оптимизации
протоколов лечения и оценки индивидуального отклика организма на терапию.

Математический фундамент системы базируется на расширенной математической модели в виде
системы обыкновенных дифференциальных уравнений, описывающей взаимодействия
между опухолевой популяцией, эффекторными клетками и концентрацией цитокинов.
Программная реализация выполнена на языке Python с использованием адаптивного
численного метода Рунге-Кутты 4(5) порядка (метод Дормана-Принса), что
обеспечило высокую точность расчетов при моделировании резких изменений
концентраций в моменты терапевтического воздействия. Разработанное
интерактивное приложение на базе PyQt5 позволяет в режиме реального времени
варьировать параметры системы и визуализировать результаты в виде расчетных
кривых и тепловых карт.

Проведенный параметрический анализ выявил доминирующую роль скорости
рекрутирования иммунных клеток (параметр $d$) как бифуркационного регулятора,
определяющего переход системы из режима автоколебаний к устойчивой ремиссии.
Внедрение интегральной метрики $AUC$ и расчет порога выживаемости ($AUC_{limit}
= 2.58 \cdot 10^{10}$ кл$\cdot$сут/мл) позволили перевести абстрактные
модельные величины в плоскость клинически значимых прогнозов.

Вычислительные эксперименты подтвердили ограниченность монотерапии цитокинами
из-за их быстрой деградации и доказали высокую эффективность комбинированных
стратегий, обеспечивающих синергетическое подавление опухоли даже на запущенных
стадиях за счет расширения «терапевтического окна». Ключевым итогом стала
успешная валидация модели на клинических данных Siu et al. (1986) по динамике
лимфомы BCL1. Достигнутый коэффициент детерминации $R^2 = 0.4706$ подтвердил
адекватность подхода в воспроизведении ключевых фаз роста и регрессии опухолевого очага.

Таким образом, цель работы была полностью достигнута. Созданный инструментарий
представляет научно-практическую ценность для предварительного анализа
иммунотерапевтических протоколов и имеет потенциал для дальнейшего развития
путем включения механизмов миграции клеток и пространственной неоднородности ткани.

В процессе выполнения научно-исследовательской работы были освоены следующие
компетенции, предусмотренные учебным планом.

УК-1. Способен осуществлять поиск, критический анализ и синтез информации,
применять системный подход для решения поставленных задач. В рамках работы
проводился систематический анализ научной литературы для формирования
3D-математической модели и поиска клинических данных Siu et al. для её
последующей валидации. Системный подход был применен при декомпозиции общей
цели исследования на подзадачи математического описания, программной реализации
и проведения вычислительных экспериментов.

УК-2. Способен определять круг задач в рамках поставленной цели и выбирать
оптимальные способы их решения, исходя из действующих правовых норм, имеющихся
ресурсов и ограничений. Для достижения цели исследования был сформулирован
перечень взаимосвязанных задач и обоснован выбор стека технологий (Python,
SciPy, PyQt5), являющегося оптимальным инструментом для численного решения
систем ОДУ и визуализации данных.

УК-3. Способен осуществлять социальное взаимодействие и реализовывать свою роль
в команде. В ходе выполнения НИР осуществлялось активное взаимодействие с
научным руководителем и профильными специалистами кафедры, что позволило
скорректировать биологическую интерпретацию результатов моделирования.

УК-4. Способен осуществлять деловую коммуникацию в устной и письменной формах
на государственном языке Российской Федерации и иностранном(ых) языке(ах).
Устная коммуникация реализовывалась в рамках обсуждения хода практики, а
письменная --- в процессе подготовки отчета и изучения фундаментальных
англоязычных источников Siu et al. и Kuznetsov et al.

УК-6. Способен управлять своим временем, выстраивать и реализовывать траекторию
саморазвития на основе принципов образования в течение всей жизни. Процесс
разработки ПО потребовал самостоятельного освоения методов глобальной
оптимизации (дифференциальная эволюция) и издательской системы \LaTeX{} для
оформления научной документации.

ПК-1. Способен проводить научно-исследовательские и опытно-конструкторские
разработки. В рамках работы был создан программно-модельный комплекс,
позволяющий исследовать динамику опухолевых клеток и эффекторного иммунитета
под внешним воздействием.

ПК-2. Способен использовать методы компьютерного моделирования, современные
методы обработки информации и оформлять полученные результаты в виде отчетов,
презентаций и докладов. В работе применены методы численного интегрирования
(метод Дормана-Принса) для исследования чувствительности системы и визуализации
результатов в виде динамических кривых и двумерных тепловых карт.

ПК-3. Способен осуществлять интеграцию программных модулей и компонент, а также
проводить проверку их работоспособности. Реализована интеграция вычислительного
ядра на базе SciPy с графическим интерфейсом на PyQt5, а корректность работы
системы подтверждена процедурой верификации на реальных данных с расчетом
коэффициента детерминации $R^2$.

ПК-4. Способен выполнять работы по созданию (модификации) и сопровождению
информационных систем. Разработанная информационная система представляет собой
расширяемую платформу, позволяющую добавлять новые уравнения (например,
пространственные эффекты) без изменения архитектуры взаимодействия GUI и решателя.

ПК-5. Способен анализировать требования к программному обеспечению и
разрабатывать технические спецификации на программные компоненты. На этапе
проектирования были проанализированы функциональные требования к системе
поддержки принятия решений и определены необходимые параметры взаимодействия
компонентов модели.

ПК-6. Способен проводить концептуальное, функциональное и логическое
проектирование информационных систем. Логическая структура ПО спроектирована с
использованием UML-диаграмм последовательности, что обеспечило изоляцию
математической логики от визуального представления.

ПК-7. Способен проектировать пользовательские интерфейсы. Спроектирован и
реализован многовкладочный графический интерфейс с интерактивными элементами
управления (слайдерами), обеспечивающий удобство проведения параметрического
анализа и моделирования сценариев терапии.